% ALPLA People Strategy 2027-2030
% FehrAdvice & Partners AG
% Version 2.0 - Februar 2026

\documentclass[11pt,a4paper]{article}

% Packages
\usepackage[utf8]{inputenc}
\usepackage[T1]{fontenc}
\usepackage[ngerman]{babel}
\usepackage{geometry}
\usepackage{graphicx}
\usepackage{xcolor}
\usepackage{array}
\usepackage{booktabs}
\usepackage{fancyhdr}
\usepackage{titlesec}
\usepackage{enumitem}
\usepackage{hyperref}
\usepackage{tikz}
\usepackage{tcolorbox}

% FehrAdvice Corporate Design Colors
\definecolor{trustblue}{RGB}{2,64,121}
\definecolor{trustbluelight}{RGB}{212,228,244}
\definecolor{fehrgray}{RGB}{37,33,42}
\definecolor{fehrlightgray}{RGB}{245,245,245}

% Page geometry
\geometry{
    top=25mm,
    bottom=25mm,
    left=25mm,
    right=25mm
}

% Headers and footers
\pagestyle{fancy}
\fancyhf{}
\fancyhead[L]{\textcolor{gray}{\small ALPLA People Strategy}}
\fancyhead[R]{\textcolor{gray}{\small FehrAdvice \& Partners AG}}
\fancyfoot[C]{\textcolor{gray}{\small Vertraulich / Confidential}}
\fancyfoot[R]{\textcolor{gray}{\small Seite \thepage}}
\renewcommand{\headrulewidth}{0.4pt}
\renewcommand{\footrulewidth}{0.4pt}

% Section formatting
\titleformat{\section}{\Large\bfseries\color{trustblue}}{\thesection}{1em}{}
\titleformat{\subsection}{\large\bfseries\color{trustblue}}{\thesubsection}{1em}{}
\titleformat{\subsubsection}{\normalsize\bfseries\color{trustblue}}{\thesubsubsection}{1em}{}

% Hyperref setup
\hypersetup{
    colorlinks=true,
    linkcolor=trustblue,
    urlcolor=trustblue,
    pdfauthor={FehrAdvice \& Partners AG},
    pdftitle={ALPLA People Strategy 2027-2030}
}

% Custom environments
\newtcolorbox{recbox}{
    colback=trustbluelight,
    colframe=trustblue,
    boxrule=0pt,
    leftrule=4pt,
    arc=0pt,
    outer arc=0pt
}

\newtcolorbox{whitebox}{
    colback=white,
    colframe=gray,
    boxrule=0pt,
    leftrule=4pt,
    arc=0pt,
    outer arc=0pt
}

\begin{document}

% ===========================================
% TITLE PAGE
% ===========================================
\thispagestyle{empty}
\begin{tikzpicture}[remember picture, overlay]
    \fill[trustblue] (current page.north west) rectangle ([yshift=-30mm]current page.north east);
    \fill[trustblue] (current page.south west) rectangle ([yshift=20mm]current page.south east);
\end{tikzpicture}

\vspace*{30mm}

\begin{center}
    % Logo placeholder
    \includegraphics[width=50mm]{image_1.png}

    \vspace{25mm}

    {\color{trustblue}\small\MakeUppercase{PROPOSAL}}

    \vspace{10mm}

    {\color{trustblue}\Huge\bfseries People Strategy\\2027--2030}

    \vspace{5mm}

    {\color{trustblue}\large Begleitung einer evidenzbasierten HR-Strategieentwicklung}

    \vspace{40mm}

    {\color{trustblue}\small Erstellt für}

    \vspace{2mm}

    {\large Andrea Becker, Head of Corporate People Development, ALPLA Group}

    \vspace{5mm}

    {\color{fehrgray} Von: FehrAdvice \& Partners AG}

    \vspace{20mm}

    {\color{gray} Februar 2026 $\cdot$ Version 2.0}
\end{center}

\vfill

\begin{tikzpicture}[remember picture, overlay]
    \node[anchor=south, text=white] at ([yshift=10mm]current page.south) {Vertraulich / Confidential};
\end{tikzpicture}

\newpage

% ===========================================
% EXECUTIVE SUMMARY
% ===========================================
\section*{Executive Summary}
\addcontentsline{toc}{section}{Executive Summary}

ALPLA steht vor der strategischen Herausforderung, eine kohärente \textbf{People Strategy} zu entwickeln, die auf die Gesamtstrategie einzahlt. In einem zunehmend kompetitiven Arbeitsmarkt erfordert dies einen systematischen, evidenzbasierten Ansatz, der drei Kernbereiche integriert:

\subsubsection*{Die drei Säulen der People Strategy:}
\begin{enumerate}
    \item \textbf{Recruiting-Strategie} -- Wie gewinnen wir die richtigen Fachkräfte?
    \item \textbf{Talent Management} -- Wie entwickeln wir Mitarbeitende systematisch?
    \item \textbf{Retention-Strategie} -- Wie binden wir kritisches Wissen langfristig?
\end{enumerate}

\textbf{Unser Ansatz:}\\
Wir begleiten ALPLA bei der Ausarbeitung der HR-Strategie auf Basis der ALPLA Gesamtstrategie. Durch eine Kombination von \textbf{klassischen Analysemethoden} (z.B. SWOT, Porter's 5 Forces) und \textbf{verhaltensökonomischen Perspektiven} entwickeln wir Strategien, die nicht nur konzeptionell überzeugend, sondern auch implementierbar sind.

\textbf{Methodik:}\\
Wir arbeiten mit \textbf{agilem Projektmanagement in Sprints}, bei dem iterativ die Analyse entwickelt und Zwischenergebnisse gesounded werden.

\begin{recbox}
\textbf{Unsere Empfehlung:}\\
Modell M (Strategieentwicklung) über 3--4 Monate mit 30 Mini-Workshops, das alle Teilstrategien vollständig entwickelt und zu einer integrierten Gesamtstrategie zusammenführt.
\end{recbox}

\subsection*{Inhaltsverzeichnis}
\begin{enumerate}[label=\arabic*.]
    \item Ausgangslage und Herausforderung
    \item Zielsetzung
    \item Geplantes Vorgehen
    \item Unser Leistungsangebot (Modelle S, M, L)
    \item Leistungsübersicht und Mehrwert
\end{enumerate}

\newpage

% ===========================================
% 1. AUSGANGSLAGE
% ===========================================
\section{Ausgangslage und Herausforderung}

\subsection{Strategischer Kontext}

Die ALPLA Gesamtstrategie definiert ambitionierte Wachstumsziele. Um diese zu erreichen, benötigt ALPLA eine HR-Strategie, die den Anforderungen des Marktes und der Organisation gleichermassen gerecht wird.

\subsection{Die zentralen Herausforderungen}

\begin{table}[h]
\centering
\begin{tabular}{p{3cm}p{5cm}p{5cm}}
\toprule
\textbf{Herausforderung} & \textbf{Beschreibung} & \textbf{Implikation} \\
\midrule
\textbf{Fachkräftemangel} & Intensiver Wettbewerb um qualifizierte Mitarbeitende & Notwendigkeit eines differenzierten Arbeitgeberangebots \\
\textbf{Alternde Belegschaft} & Kritisches Wissen droht verloren zu gehen & Systematischer Wissenstransfer und Nachfolgeplanung \\
\textbf{Arbeitgeberattraktivität} & Unklare Positionierung im Wettbewerb & Definition eines «Rahmenpakets» für Attraktivität \\
\textbf{Mitarbeiterbindung} & Ungewollte Fluktuation bindet Ressourcen & Frühwarnsystem und proaktive Bindungsstrategien \\
\textbf{Kompetenzentwicklung} & Skill-Gaps bei strategisch wichtigen Positionen & Strukturiertes Talent Management \\
\bottomrule
\end{tabular}
\end{table}

\section{Zielsetzung}

\begin{whitebox}
\textbf{Primärziel:} Entwicklung einer kohärenten, evidenzbasierten \textbf{People Strategy 2027--2030}, die auf die ALPLA Gesamtstrategie einzahlt und sowohl dem Wettbewerbsumfeld als auch dem Organisationskontext gerecht wird.
\end{whitebox}

\subsubsection*{Teilziele}
\begin{itemize}
    \item \textbf{Recruiting-Strategie:} Definition Rahmenpaket «Attraktiver Arbeitgeber», Wettbewerbsanalyse, Massnahmenplan Time-to-Hire.
    \item \textbf{Talent Management-Strategie:} Strategie zur Identifikation kritischer Kompetenzen, Design von Entwicklungsprogrammen, Nachfolgeplanung.
    \item \textbf{Retention-Strategie:} Frühwarnsystem für Mitarbeiterzufriedenheit, Konzept für Wissenstransfer.
    \item \textbf{Gesamtstrategie:} Integration aller Teilstrategien unter Berücksichtigung ihrer Komplementaritäten, Priorisierung auf Basis von Wettbewerbsumfeld und Organisationskontext, Implementierungs-Roadmap.
\end{itemize}

\newpage

% ===========================================
% 3. GEPLANTES VORGEHEN
% ===========================================
\section{Geplantes Vorgehen}

\subsection{Methodischer Ansatz}

Wir arbeiten mit \textbf{agilem Projektmanagement in Sprints}, bei dem iterativ die Analyse entwickelt und Zwischenergebnisse gesounded werden:

\begin{center}
\textbf{Analyse} $\rightarrow$ \textbf{Hypothesen} $\rightarrow$ \textbf{Validierung} $\rightarrow$ \textbf{Sounding} $\rightarrow$ \textbf{Iteration} $\circlearrowleft$
\end{center}

\subsection{Informationsgewinnung}

\begin{table}[h]
\centering
\begin{tabular}{lll}
\toprule
\textbf{Typ} & \textbf{Methode} & \textbf{Beispiele} \\
\midrule
\textbf{Extern} & Klassische Analyse & Marktanalyse, Wettbewerbsanalyse, Porter's 5 Forces, Benchmarks \\
\textbf{Intern} & Agile Workshops & Mini-Workshops, Stakeholder-Interviews, HR-Datenauswertung \\
\bottomrule
\end{tabular}
\end{table}

\subsection{Der Strategieprozess in 4 Phasen}

\begin{table}[h]
\centering
\small
\begin{tabular}{p{3.2cm}p{3.2cm}p{3.2cm}p{3.2cm}}
\toprule
\textbf{PHASE 1} & \textbf{PHASE 2} & \textbf{PHASE 3} & \textbf{PHASE 4} \\
\textbf{Recruiting} & \textbf{Talent Mgmt} & \textbf{Retention} & \textbf{Gesamtstrategie} \\
\midrule
$\bullet$ Marktanalyse & $\bullet$ Status-Quo-Analyse & $\bullet$ Marktanalyse & $\bullet$ Integration \\
$\bullet$ Wettbewerbsanalyse & $\bullet$ Potenzialanalyse & $\bullet$ Status-Quo-Analyse & $\bullet$ Konsistenzprüfung \\
$\bullet$ Status-Quo-Analyse & $\bullet$ Ressourcenallokation & $\bullet$ MA-Umfrage Konzept & $\bullet$ Priorisierung \\
$\bullet$ Potenzialanalyse & & $\bullet$ Potenzialanalyse & $\bullet$ Governance \\
\midrule
\multicolumn{4}{c}{\cellcolor{trustbluelight} $\blacktriangledown$ Meilenstein 1 \hfill $\blacktriangledown$ Meilenstein 2 \hfill $\blacktriangledown$ Meilenstein 3 \hfill $\blacktriangledown$ Abschluss} \\
\bottomrule
\end{tabular}
\end{table}

\subsection{Verhaltensökonomische Perspektive}

Alle Analysen erfolgen nicht nur klassisch, sondern zusätzlich unter Berücksichtigung verhaltensrelevanter Merkmale:
\begin{itemize}
    \item \textbf{Marktanalyse:} + Entscheidungsverhalten von Kandidaten
    \item \textbf{Wettbewerbsanalyse:} + Wahrgenommene vs. tatsächliche Attraktivität
    \item \textbf{Status-Quo-Analyse:} + Verhaltensbarrieren und -treiber
\end{itemize}

\newpage

% ===========================================
% 4. LEISTUNGSANGEBOT
% ===========================================
\section{Unser Leistungsangebot}

\subsection{Modell S -- Strategie-Review}

\textbf{Ziel:} Gesamtbetrachtung der HR-Strategie und Health-Check.

\begin{itemize}
    \item \textbf{Leistungen:} 5 Mini-Workshops, SWOT-Schnellanalyse, Priorisierte Handlungsempfehlungen.
    \item \textbf{Lieferobjekte:} Strategie-Assessment, Top 10 Handlungsempfehlungen, Management-Präsentation.
\end{itemize}

\textbf{Parameter:} 4--6 Wochen | EUR 20'000 | 1--2 Consultants

\vspace{5mm}
\hrule
\vspace{5mm}

\begin{recbox}
\subsection{Modell M -- Strategieentwicklung $\star$ EMPFEHLUNG}

\textbf{Ziel:} Vollständige Entwicklung aller drei Teilstrategien und Integration zur Gesamtstrategie inkl. Roadmap.
\end{recbox}

\subsubsection*{Leistungen (Übergreifend)}
\begin{itemize}
    \item \textbf{30 Mini-Workshops} verteilt auf 4 Phasen.
    \item Vollständige interne und externe Analyse pro Teilbereich inkl. verhaltensökonomischer Perspektive.
    \item Iteratives Sounding nach jeder Phase (4 Meilensteine).
\end{itemize}

\subsubsection*{Phasen \& Inhalte}
\begin{itemize}
    \item \textbf{Phase 1 (Recruiting):} Markt-/Wettbewerbsanalyse, Status-Quo, Potenzialanalyse.
    \item \textbf{Phase 2 (Talent Mgmt):} Bedarfsanalyse (Kompetenzen 2027--2030), Nachfolgeplanung.
    \item \textbf{Phase 3 (Retention):} Fluktuationstrends, Optimierung MA-Umfrage, Fokus alternde Belegschaft.
    \item \textbf{Phase 4 (Gesamtstrategie):} Synthese, Konsistenzprüfung, Roadmap, KPI-Framework.
\end{itemize}

\subsubsection*{Lieferobjekte}
\begin{itemize}
    \item Detaillierte Strategien für Recruiting, Talent Mgmt \& Retention.
    \item Integrierte People Strategy 2027--2030 \& Implementierungs-Roadmap.
    \item Umfassende Management-Präsentation.
\end{itemize}

\begin{table}[h]
\centering
\begin{tabular}{ll}
\toprule
\textbf{Parameter} & \textbf{Wert} \\
\midrule
Dauer & 3--4 Monate \\
Investition & EUR 140'000 \\
Team & 2--3 Consultants \\
\bottomrule
\end{tabular}
\end{table}

\newpage

\subsection{Modell L -- Strategieverankerung}

\textbf{Ziel:} Alles aus Modell M, plus verhaltensbasierte Verankerung und Implementierungsbegleitung.

\subsubsection*{Zusätzliche Leistungen}
\begin{itemize}
    \item \textbf{50 Mini-Workshops} gesamt.
    \item \textbf{Change Management:} Stakeholder-Mapping, Kommunikation, Widerstandsmanagement.
    \item \textbf{Verhaltensbasierte Verankerung:} Nudging-Konzepte, Choice Architecture.
    \item \textbf{Pilotierung} ausgewählter Massnahmen \& Erfolgsmessung (3 Quartals-Reviews).
\end{itemize}

\textbf{Parameter:} 6--9 Monate | EUR 250'000 | 3--4 Consultants

% ===========================================
% 5. LEISTUNGSÜBERSICHT
% ===========================================
\section{Leistungsübersicht und Mehrwert}

\subsection{Leistungsvergleich}

\begin{table}[h]
\centering
\small
\begin{tabular}{lcccc}
\toprule
\textbf{Leistung} & \textbf{Modell S} & \textbf{Modell M} & \textbf{Modell L} \\
\midrule
Strategischer Health-Check & $\checkmark$ & $\checkmark$ & $\checkmark$ \\
Teilstrategien (Recr/Talent/Ret) & Empfehlungen & Vollständig & + Implementierung \\
Gesamtstrategie & Assessment & $\checkmark$ & $\checkmark$ \\
Implementierungs-Roadmap & -- & $\checkmark$ & $\checkmark$ \\
Change Mgmt \& Pilotierung & -- & -- & $\checkmark$ \\
\midrule
\textbf{Investition} & \textbf{EUR 19'250} & \textbf{EUR 140'250} & \textbf{EUR 250'000} \\
\bottomrule
\end{tabular}
\end{table}

\subsection{Quantifizierbarer Nutzen (Schätzung)}

\begin{table}[h]
\centering
\begin{tabular}{lll}
\toprule
\textbf{Kennzahl} & \textbf{Verbesserung} & \textbf{ALPLA Potenzial} \\
\midrule
Reduktion Fluktuation & -15\% bis -25\% & EUR 2--4 Mio. / Jahr \\
Reduktion Recruitingkosten & -10\% bis -20\% & EUR 0.5--1 Mio. / Jahr \\
\bottomrule
\end{tabular}
\end{table}

\newpage

\subsection{Warum FehrAdvice?}

\begin{table}[h]
\centering
\begin{tabular}{p{4cm}p{9cm}}
\toprule
\textbf{Merkmal} & \textbf{Beschreibung} \\
\midrule
Wissenschaftlichkeit & Prof. Ernst Fehr als wissenschaftlicher Beirat \\
Verhaltensökonomie & Einzigartige Kombination von Strategie und Behavioral Science \\
Praxiserfahrung & 200+ HR-Transformationsprojekte \\
Proprietäre Tools & BEATRIX-Suite für datengestützte Analyse \\
\bottomrule
\end{tabular}
\end{table}

\subsection{Unsere Empfehlung}

Für ALPLA empfehlen wir \textbf{Modell M (Strategieentwicklung)} als optimale Balance zwischen Tiefe und Effizienz:

\begin{itemize}
    \item[$\checkmark$] Vollständige Entwicklung aller drei Teilstrategien
    \item[$\checkmark$] Integration zu einer kohärenten Gesamtstrategie
    \item[$\checkmark$] Konkrete Implementierungs-Roadmap
    \item[$\checkmark$] Angemessener Zeit- und Ressourceneinsatz
\end{itemize}

Bei Bedarf kann das Projekt später um Implementierungsbegleitung (Elemente aus Modell L) erweitert werden.

\newpage

% ===========================================
% APPENDIX - WER WIR SIND
% ===========================================
\section*{Appendix -- Wer wir sind}
\addcontentsline{toc}{section}{Appendix -- Wer wir sind}

\textbf{FehrAdvice \& Partners AG} ist die führende \textbf{Strategie- und Managementberatung für Verhaltensökonomie} im deutschsprachigen Raum. Seit der Gründung durch Prof. Ernst Fehr (Universität Zürich) und Gerhard Fehr im Jahr 2010 begleiten wir Unternehmen, Behörden und Organisationen bei der Strategieentwicklung und Gestaltung wirksamer, strategiebasierter Management- und Steuerungslogiken -- fundiert, datenbasiert und anschlussfähig an Regulatorik.

\textbf{Das methodische Fundament unserer Arbeit bildet das Behavioral Change Model (BCM):} Dieses von FehrAdvice entwickelte Modell basiert auf über 40 Jahren Verhaltensökonomie und mehr als 1.500 wissenschaftlichen Studien -- unter ständiger wissenschaftlicher Leitung von Univ. Prof. Ernst Fehr, einem der einflussreichsten und meistzitierten Ökonomen weltweit.

\subsection*{Unsere Core-Practices}

\begin{itemize}
    \item \textbf{Management-Strategie:} Evidenzbasierte Strategieentwicklung und Umsetzung
    \item \textbf{Stakeholder- und Akzeptanzmanagement:} Systematisches Stakeholdermapping und Reaktanzanalysen
    \item \textbf{Kundenstrategie:} Datenbasierte Segmentierung und gezielte Stakeholderansprache
    \item \textbf{Governance \& Compliance:} Governance-Architekturen, Regulatorik, Monitoring
    \item \textbf{Digitale Transformation:} Entwicklung digitaler Entscheidungsarchitekturen (BEATRIX)
\end{itemize}

\begin{recbox}
\textbf{Kurz gesagt:}\\
FehrAdvice macht Strategie und Veränderung gestaltbar -- evidenzbasiert, praxisnah und zukunftssicher. Nicht durch Kontrollillusion, sondern durch bessere strategische und operative Entscheidungen.
\end{recbox}

\vfill

\begin{center}
\hrule
\vspace{5mm}
\textbf{FehrAdvice \& Partners AG}\\
Binzmühlestrasse 170A $\cdot$ 8050 Zürich $\cdot$ Schweiz\\
\href{mailto:info@fehradvice.com}{info@fehradvice.com} $\cdot$ +41 44 256 79 00 $\cdot$ \href{https://www.fehradvice.com}{www.fehradvice.com}
\vspace{5mm}
\hrule
\vspace{3mm}
{\small\color{gray} © 2026 FehrAdvice \& Partners AG. Alle Rechte vorbehalten.}
\end{center}

\end{document}
